\section{Introduction}
\label{sec:intro}

Numerical weather prediction (NWP) output is systematically biased at any
single station or grid point: model resolution, physical
parameterizations, and terrain representation err in a station-specific
and often persistent way, and statistical postprocessing exists
specifically to correct this systematic component rather than relying on
raw ensemble output directly \citep{vannitsem_2020_postprocessing_review,gneiting_2005_emos}.
Beyond Ensemble Model Output Statistics (EMOS)'s original Gaussian
regression \citep{gneiting_2005_emos}, the field has extended
postprocessing along several axes: mixture and non-Gaussian predictive
families \citep{baran_2015_mixture_emos}, target-specific reformulations
for variables EMOS's Gaussian assumption fits poorly
\citep{scheuerer_2013_precip_emos}, explicit calibration-enforcing
extensions \citep{wilks_2018_enforcing_calibration}, and
permutation-invariant neural architectures that treat ensemble members as
an unordered set \citep{hoehlein_2023_permutation_invariant}. The
Distributional Regression Network (DRN, \citealp{rasp_2018_drn}) is this
paper's own neural baseline (Section~\ref{sec:methods}).

Every one of these methods shares the same requirement: a multi-year
archive of paired ensemble/observation history at the specific station
being corrected. \citet{vannitsem_2020_postprocessing_review}'s review of
the field flags this as a structural limitation, not an incidental one --
and it is exactly what a newly commissioned station, a station following
an NWP model version upgrade that breaks the archive's stationarity, or a
station in a sparsely observed region cannot supply. Trained
postprocessing's accuracy gains, in other words, come with an entry cost
that a purely zero-shot alternative would not.

Time series foundation models (TSFMs) -- large, pretrained sequence models
applied without any task-specific training or parameter updates on the
target series -- offer a structurally different answer: no
station-specific training archive at all. This is a different proposition
from the spatial, full-field AI weather-prediction models that have
separately transformed NWP itself
\citep{lam_2023_graphcast,price_2024_gencast,bodnar_2025_aurora}: those
models replace the numerical forecast; a TSFM used for postprocessing
takes an existing ensemble forecast as input and corrects it at the
station level, one series at a time, with no spatial field anywhere in its
architecture. TimesFM-3 -- the 3.0 release of Das et al.'s decoder-only
TimesFM architecture \citep{das_2024_timesfm_icml} -- is one such model, and
its past-future covariate support (a version-specific capability sourced
directly in Section~\ref{sec:covariates}, not part of the original
architecture paper) is the specific capability this paper exploits: the
ensemble forecast can be supplied as a covariate known across the entire
forecast horizon, not just in the past, letting a zero-shot model use NWP
guidance the way a postprocessing method would, without training on this
benchmark's data at all.

To our knowledge, no existing work directly measures the training-data
size at which a station's own trained postprocessing model overtakes a
zero-shot time series foundation model given the same NWP guidance -- the
question this paper answers. On the identical EUPPBench benchmark this
paper uses, \citet{vanpoecke_2024_transformer},
\citet{mlakar_2023_spline_flows}, and \citet{landry_2025_fmap} each
propose a trained neural or flow-based postprocessing architecture, and
\citet{allen_2023_multivariate_calibration} evaluates multivariate
calibration on it -- none compares against a zero-shot foundation model.
\citet{baran_2026_sharpness_penalty} studies the same sharpness axis this
paper's calibration-sharpness result depends on, on the same benchmark,
but not the zero-shot-vs-trained breakpoint question.
\citet{sun_2026_streamflow_zeroshot} pose a structurally identical
question -- is a zero-shot time series foundation model competitive with
domain-trained methods -- in hydrology, not weather postprocessing;
\citet{bharadwaj_2026_air_quality_arena} benchmark TSFM zero-shot skill
across regions in air quality forecasting, a related but distinct target
variable and evaluation design. \citet{wessel_2024_leadtime_continuous}
and \citet{muschinski_2023_mos_random_forests} are the closest
small-training-data postprocessing rivals to this paper's own breakpoint
curve, targeting the same low-data regime with different
(lead-time-continuous, random-forest) trained architectures rather than a
zero-shot comparison. None of these measures where the crossing point
actually is.

This paper's contribution is threefold. (1) At full training data, every
trained method (EMOS pooled, EMOS local, DRN) beats TimesFM-3's
Continuous Ranked Probability Score (CRPS, Section~\ref{sec:results-full}).
A one-parameter variance-inflation
rescaling of the raw ensemble does too: a fixed-climatological variant
beats it with no training data at all, and the CRPS-optimal variant beats
it at every training size we tested it at, from \InflASmallestNDays{}
days upward (Section~\ref{sec:results-full}) -- on CRPS, the answer to
this paper's central question is negative for TimesFM-3, and not
narrowly so. This advantage starts at very small training sizes: pooled
EMOS's CRPS crosses TimesFM-3's at \BpCrpsPooledCases{} cases
(\BpCrpsPooledDays{} calendar days, Section~\ref{sec:results-breakpoint}),
but that raw crossing is not itself statistically significant; the
advantage is significant from \BpCrpsPooledSigK{} cases
(\BpCrpsPooledSigDays{} calendar days) onward under day-blocking, our
primary block definition, for both EMOS pooled and EMOS local alike --
the station-blocked \(t\)-test does not reach significance anywhere in
the tested \(k=1\)--9 range, while the station-blocked Wilcoxon test does
from \(k=\BpCrpsPooledStationWilcoxonFirstSigK{}\) onward, the same
test-statistic divergence already reported at \SigTestNCases{} cases
(Section~\ref{sec:results-significance}), now shown to persist across the
whole low-\(k\) range. Coverage crosses more slowly
(ratios computed from the unrounded breakpoint estimates, not the
values as displayed): pooled EMOS's coverage crosses TimesFM-3's roughly
\CoverageCrpsRatioPooled{}x later than its CRPS does (\BpCoveragePooledCases{}
cases / \BpCoveragePooledDays{} days), and local EMOS's about
\CoverageCrpsRatioLocal{}x later (\BpCoverageLocalCases{} cases /
\BpCoverageLocalDays{} days) -- so "a few days is enough" holds for CRPS
specifically, not for calibration. These comparisons are between methods
given different information: TimesFM-3's covariate channel received the
ensemble mean but not ensemble spread in this study, while EMOS and DRN
received spread but no observation history (Section~\ref{sec:covariates},
Table~\ref{tab:info-channels}); supplying spread as a second covariate
does not close this calibration-sharpness gap either
(Section~\ref{sec:discussion-covariate-check}). (2) On the
calibration-sharpness axis, TimesFM-3 holds a real advantage over a
trivial rescaling: even a leakage-optimistic multiplier fit directly on
test-set coverage -- an upper bound no legitimately deployable trivial
rescaling could exceed -- needs \InflLeakyWidth{} K, still wider than
TimesFM-3's \WidthTsfmFull{} K (Section~\ref{sec:discussion-trivial}).
(3) TimesFM-3's short-lead CRPS advantage first
flips sign at \FirstSignFlipHours{}~h and reverses durably by
\DurableCrossoverHours{}~h (Section~\ref{sec:results-leadtime}); the
mechanism behind that advantage has two independent parts. The
observation-history part was tested directly and does not support an
information-asymmetry explanation -- a persistence-augmented EMOS
closes only \PersistenceGapClosedPct{}\% of the gap -- but the test
itself cannot cleanly distinguish an information-set explanation from an
architectural one (Section~\ref{sec:discussion-mechanism}). The
ensemble-spread/calibration part was also tested directly, by supplying
TimesFM-3 with ensemble spread as a second covariate: it worsens CRPS
and shifts the model along the calibration-sharpness trade-off without
closing the calibration deficit, even though the channel demonstrably
accepts and uses the input (Section~\ref{sec:discussion-covariate-check}).
Together these results argue that a frozen, zero-shot time series
foundation model is, for this task, mostly not competitive with a few
days of station-specific training data on CRPS -- with one genuine
exception that survives even a leakage-optimistic trivial baseline (the
calibration-sharpness axis), and a short-lead advantage that the
information-asymmetry explanation does not account for. The evidence
presented here is for one
time series foundation model, TimesFM-3; whether it generalizes to other
pretrained time series models (e.g.\ Chronos-2, Moirai-2) is untested.
